\documentclass[12pt,a4paper]{article}
\usepackage[utf8]{inputenc}
\usepackage[T1]{fontenc}
\usepackage{amsmath,amssymb,amsthm}
\usepackage{geometry}
\usepackage{hyperref}
\usepackage{booktabs}

\geometry{margin=1in}

\newtheorem{theorem}{Theorem}
\newtheorem{axiom}{Axiom}
\newtheorem{proposition}{Proposition}
\newtheorem{lemma}{Lemma}
\newtheorem{conjecture}{Conjecture}

\theoremstyle{definition}
\newtheorem{definition}{Definition}

\theoremstyle{remark}
\newtheorem{remark}{Remark}

\title{\textbf{A self-consistency hierarchy for timeless quantum gravity}\\[6pt]
{\large DRAFT}}

\author{Will Regelmann\thanks{Corresponding author.} \quad Claude (Anthropic)\thanks{AI
assistant; contributed to mathematical development and analysis.}}

\date{March 2026}

\begin{document}
\maketitle

\begin{abstract}
We propose a novel formulation of quantum gravity consistent with a block
universe interpretation: no fundamental Hilbert space, no evolution parameter,
no background structure. The proposal is organized as a hierarchy of
self-consistency constraints at four levels of coarse-graining: topological
(piecewise-linear 4-manifolds with intersection form constraint), smooth
(exotic smooth structures on the topological type), metric (the semiclassical
Einstein equation as a fixed point, studied in the companion
paper~\cite{fixed_point}), and effective quantum mechanical (density matrices
as marginals of the smooth-structure measure). The Hilbert space of standard
quantum mechanics is not postulated---it is the derived structure of the
marginal measure restricted to a subregion of the block. The four-dimensionality
of the universe is not an input: four is the unique dimension where smooth
structures are rich enough (uncountably many inequivalent structures on
$\mathbb{R}^4$, versus exactly one in all other dimensions) to support the
required self-consistency structure. The Banach contraction result of the
companion paper~\cite{fixed_point} anchors the perturbative regime at the
metric level. We identify four priority open problems, of which the bridge
from intersection form signature to metric signature is the most critical.

\bigskip\noindent\textit{Note:} This paper is a companion to ``Fixed-point
existence of self-consistent semiclassical gravity'' and ``Gaussian temporal
profile of gravitational decoherence from the Einstein--Langevin equation.''
It proposes the quantum-gravitational framework within which the results of
those papers are understood as the metric-level incarnation of a more
fundamental topological self-consistency structure.
\end{abstract}

%======================================================================
\section{Introduction}
\label{sec:intro}
%======================================================================

The foundational tension in quantum gravity is well known: general relativity
requires no preferred time parameter, while standard quantum mechanics requires
one. The Wheeler--DeWitt equation $\hat{H}\Psi = 0$~\cite{dewitt,wheeler} is
the canonical attempt at resolution---the Hamiltonian constraint makes the
wave functional timeless. But every known formulation of $\hat{H}\Psi = 0$
introduces hidden temporal or foliation structure at some stage of its
derivation~\cite{kuchar}. The canonical ADM approach requires a $3+1$
decomposition; the covariant path integral hides foliation dependence in the
Wick rotation, boundary conditions, and measure ambiguity; the WKB
semiclassical limit inserts time via Born--Oppenheimer
adiabaticity~\cite{halliwell,banks}.

The difficulty runs deeper than any specific implementation. Every existing
formulation of quantum gravity inherits Hilbert space from quantum mechanics,
and Hilbert space carries implicit temporal structure: states are vectors that
``evolve,'' inner products require a time-slice to define, and the Born rule
describes what happens ``when you measure.'' Taking the block universe
requirement seriously---the universe is a four-dimensional geometric structure
satisfying constraints, not a state evolving in time---forces a more radical
conclusion: \emph{the Hilbert space cannot be fundamental}. It must emerge as
an effective description for subsystems with limited access to the full block.

This paper develops the formulation that follows from this conclusion. The
organizing principle is:

\begin{axiom}[No evolution parameter]
\label{ax:timeless}
The fundamental description contains no evolution parameter. All equations
are constraints on the structure of the four-dimensional block, not rules for
generating it. A formulation satisfies this axiom if and only if it can be
stated entirely as a constraint on a four-dimensional geometric structure with
no preferred slicing.
\end{axiom}

\begin{axiom}[Self-consistency]
\label{ax:selfconsistency}
The coarse-grained structure of the universe is the unique solution to its
own consistency constraints. The solution has no closed-form expression but
is the convergent limit of an iterative process~(Conjecture). Multiple
fine-grained configurations are compatible with the same coarse-grained
solution~(Conjecture).
\end{axiom}

\begin{axiom}[Geometric purity]
\label{ax:geometric}
The consistency constraints are purely geometric. No non-geometric entities
appear in the fundamental description. At the deepest level, the constraints
are topological: metric structure and matter both emerge from topology~(Conjecture).
\end{axiom}

The formulation we propose satisfies Axiom~\ref{ax:timeless} at every level.
It is organized as a hierarchy of self-consistency constraints
(Axiom~\ref{ax:selfconsistency}) that are purely geometric
(Axiom~\ref{ax:geometric}).

\medskip

\noindent\textbf{What this paper is and is not.} This is a framework proposal,
not a completed theory. Each level of the hierarchy rests on specific
mathematical conjectures, which we label carefully. The metric level has the
most rigorous support---it is the setting of the companion
paper~\cite{fixed_point}. The topological and smooth-structure levels are
research programs with clear mathematical agendas. We do not claim to have
derived the Standard Model or explained consciousness. We claim to have
identified the correct \emph{mathematical home} for the axioms, and to have
made explicit what is required to fill it in.

%======================================================================
\section{Why Hilbert space cannot be fundamental}
\label{sec:nohilbert}
%======================================================================

We record the argument that motivates abandoning Hilbert space as a
fundamental object, since it drives the entire proposal.

A Hilbert space $\mathcal{H}$ with inner product $\langle\cdot|\cdot\rangle$
carries temporal structure in three ways, any one of which is fatal for
Axiom~\ref{ax:timeless}:

\begin{enumerate}

\item \textbf{The inner product requires a time-slice.} In quantum field
theory, $\langle\Psi_1|\Psi_2\rangle = \int \Psi_1^*[\phi]\Psi_2[\phi]
\,\mathcal{D}\phi$ integrates over field configurations $\phi$ on a spatial
hypersurface at a fixed time. Defining the inner product requires choosing that
hypersurface. Different choices give inequivalent inner products and
inequivalent quantum theories.

\item \textbf{Unitary evolution requires a time parameter.} The statement
$|\Psi(t_2)\rangle = U(t_2, t_1)|\Psi(t_1)\rangle$ requires a preferred
parameter $t$. The Schr\"odinger equation $i\hbar\partial_t|\Psi\rangle =
\hat{H}|\Psi\rangle$ is defined with respect to $t$. Even in the
``frozen'' Wheeler--DeWitt formulation, the Hilbert space inner product
is defined on a spatial slice, smuggling in the foliation through the
back door.

\item \textbf{The Born rule requires a preferred decomposition into outcomes.}
The probability $|\langle\phi|\Psi\rangle|^2$ is defined relative to a
measurement basis $\{|\phi\rangle\}$. Selecting this basis is an
observer-dependent choice that the fundamental description should not require.

\end{enumerate}

\begin{proposition}[No Axiom-\ref{ax:timeless}-compliant Hilbert space formulation]
\emph{(Rigorous for (1) and (2); Sketch for (3).)}
\label{prop:nohilbert}
Any formulation of quantum gravity that takes a Hilbert space with inner product
as a fundamental input violates Axiom~\ref{ax:timeless}, because the inner
product requires a preferred spatial hypersurface (point (1)) or a preferred
time parameter (point (2)).
\end{proposition}

\noindent\emph{Argument.} Point (1) follows from the standard QFT definition of
the Fock space inner product. The functional integral $\int\mathcal{D}\phi$ is
over field configurations on a Cauchy surface. Different Cauchy surfaces give
different inner products (Reeh-Schlieder theorem, Unruh effect). For
Axiom~\ref{ax:timeless} compliance, no Cauchy surface can be preferred.
Point (2) is immediate from the definition of unitary evolution. \hfill$\square$

\medskip

The conclusion is that a formulation satisfying Axiom~\ref{ax:timeless} must
not take Hilbert space as input. The Hilbert space, if it appears at all, must
be derived---a structure that emerges from the fundamental description for
subsystems with limited access to the block.

%======================================================================
\section{The self-consistency hierarchy}
\label{sec:hierarchy}
%======================================================================

We propose a four-level hierarchy of self-consistency constraints. The levels
are organized by \emph{coarse-graining}: Level~0 is the coarsest description,
Level~3 the finest accessible to a subsystem. The levels are not temporal
stages. They are simultaneous aspects of the same four-dimensional block,
described at different resolutions. The constraint operates simultaneously at
all levels: a valid block must satisfy the constraints at every level, and these
constraints are not independent---they interlock.

\begin{definition}[The four-level hierarchy]
\label{def:hierarchy}
\quad
\begin{itemize}
\item[\emph{Level 0.}] \emph{Topological.} The fundamental object is a
piecewise-linear (PL) 4-manifold $\mathcal{M}$, equipped with no additional
structure. The consistency constraint selects an intersection form $Q_\mathcal{M}$
on $H_2(\mathcal{M};\mathbb{Z})$.

\item[\emph{Level 1.}] \emph{Smooth.} Given the topological type of
$\mathcal{M}$, the space of compatible smooth structures $\mathrm{Sm}(\mathcal{M})$
is the configuration space of the theory. A self-consistency measure $\mu^*$ on
$\mathrm{Sm}(\mathcal{M})$ replaces both the wavefunction and the path integral
measure of standard formulations.

\item[\emph{Level 2.}] \emph{Metric.} Each smooth structure $\sigma$ constrains
the space of compatible Riemannian metrics. The semiclassical Einstein equation
$G_{\mu\nu} = 8\pi G\langle\hat{T}_{\mu\nu}\rangle_g$ is the fixed-point
condition that selects the metric at this level; its existence is established in
the companion paper~\cite{fixed_point}. The mechanism by which a smooth
structure constrains the metric---beyond the requirement of compatibility---is
not fully developed: the Yamabe problem provides a metric of constant scalar
curvature within each conformal class, and Donaldson/Seiberg--Witten invariants
distinguish smooth structures, but neither alone determines a unique metric.
The claim that a smooth structure together with the self-consistency constraint
determines a metric is a Conjecture.

\item[\emph{Level 3.}] \emph{Effective quantum mechanics.} Marginalizing the
smooth-structure measure $\mu^*$ over the complement of a subregion
$\mathcal{R}\subset\mathcal{M}$ produces an effective measure on the
configuration space of $\mathcal{R}$. This marginal is the effective density
matrix of the subsystem. The Hilbert space $\mathcal{H}_\mathcal{R}$ is the
$L^2$ space of this marginal.
\end{itemize}
\end{definition}

The remainder of this section develops each level and specifies the
self-consistency constraints that link them.

%----------------------------------------------------------------------
\subsection{Level 0: The topological constraint}
\label{sec:level0}
%----------------------------------------------------------------------

Let $\mathcal{M}$ be a compact, oriented, simply-connected PL 4-manifold.
(The simply-connected assumption is used below to invoke Freedman's
classification theorem; extensions to non-simply-connected manifolds are an
open problem.)

The intersection form $Q_\mathcal{M}: H_2(\mathcal{M};\mathbb{Z}) \times
H_2(\mathcal{M};\mathbb{Z}) \to \mathbb{Z}$ is defined by
$Q_\mathcal{M}(\alpha,\beta) = \langle \alpha \cup \beta, [\mathcal{M}]\rangle$,
the Kronecker pairing of the cup product with the fundamental class. For a
closed oriented 4-manifold, $Q_\mathcal{M}$ is symmetric and unimodular (the
latter by Poincar\'e duality).

\begin{theorem}[Freedman, 1982~\cite{freedman}]
\emph{(Rigorous.)}
\label{thm:freedman}
Every unimodular symmetric bilinear form over $\mathbb{Z}$ is realized as
the intersection form of a simply-connected compact topological
$4$-manifold. The topological type of the manifold is determined by the
intersection form and the Kirby--Siebenmann invariant.
\end{theorem}

Freedman's theorem classifies topological 4-manifolds. Donaldson's theorem
constrains which of these admit smooth structures:

\begin{theorem}[Donaldson, 1983~\cite{donaldson}]
\emph{(Rigorous.)}
\label{thm:donaldson}
If a simply-connected smooth closed $4$-manifold has definite intersection
form $Q$, then $Q$ is diagonalizable over $\mathbb{Z}$ (i.e., $Q \cong
\pm I$ in some basis). In particular, the $E_8$ form and other non-diagonalizable
definite forms are not realized as intersection forms of smooth $4$-manifolds.
\end{theorem}

The gap between Theorems~\ref{thm:freedman} and~\ref{thm:donaldson}---topological
manifolds that exist but admit no smooth structure---is where the richest physics
lives. It is also where Axiom~\ref{ax:selfconsistency} does its first work:

\begin{conjecture}[Topological self-consistency]
\label{conj:level0}
\emph{(Conjecture.)}
The consistency constraints of Axiom~\ref{ax:selfconsistency}, applied at
Level~0, select a specific unimodular form $Q^*$ as the intersection form of
$\mathcal{M}$. The constraint is: $Q^*$ must be compatible with a smooth
structure (by Donaldson's theorem, $Q^*$ must be indefinite or diagonalizable),
and the smooth structures compatible with $Q^*$ must, at Level~1, produce a
measure $\mu^*$ whose metric-level content is consistent with the matter
content that determined $Q^*$ in the first place.
\end{conjecture}

The self-referential structure of Conjecture~\ref{conj:level0} is not a defect.
It is the precise content of Axiom~\ref{ax:selfconsistency}: the universe is
the unique solution to its own consistency constraints.

%----------------------------------------------------------------------
\subsection{Level 1: The smooth-structure measure}
\label{sec:level1}
%----------------------------------------------------------------------

The most remarkable fact in four-dimensional topology is the following:

\begin{theorem}[Donaldson, 1983; Taubes, 1987~\cite{donaldson,taubes}]
\emph{(Rigorous.)}
\label{thm:exotic}
The smooth manifold $\mathbb{R}^4$ admits uncountably many mutually
non-diffeomorphic smooth structures. In every other dimension $n \neq 4$,
$\mathbb{R}^n$ admits a unique smooth structure up to diffeomorphism.
\end{theorem}

This is the mathematical fact that makes four dimensions special. The space
$\mathrm{Sm}(\mathcal{M})$ of smooth structures compatible with the PL type
of $\mathcal{M}$ is, for a four-dimensional $\mathcal{M}$, generically an
uncountable set. In any other dimension, it is a finite set (often a single
point).

\begin{definition}[The smooth-structure measure]
\emph{(Conjecture.)}
\label{def:mu}
The self-consistency measure $\mu^*$ on $\mathrm{Sm}(\mathcal{M})$ is the
unique measure satisfying the self-consistency condition
$\mathcal{F}(\mu^*) = \mu^*$, where $\mathcal{F}$ is the map defined by:
given a trial measure $\mu$ on $\mathrm{Sm}(\mathcal{M})$, compute the
induced metric and matter content at each smooth structure (via Level~2),
compute the resulting stress-energy, solve the Einstein constraint for a
new geometry, and update the weights on $\mathrm{Sm}(\mathcal{M})$
accordingly.
\end{definition}

\begin{remark}[Relation to the path integral]
The smooth-structure measure $\mu^*$ is an analogue of the gravitational path
integral, but with a crucial structural difference. The standard gravitational
path integral integrates over a space of metrics, which is a continuous
infinite-dimensional space with an undefined measure. The measure $\mu^*$
integrates over smooth structures on a fixed topological type. While
$\mathrm{Sm}(\mathcal{M})$ may be uncountable, its structure as a moduli space
of gauge-equivalence classes is better controlled than the space of all metrics.
In particular, Donaldson and Seiberg--Witten theory provide invariants of smooth
structures that can serve as weights in $\mu^*$, even if the full measure is not
yet defined.
\end{remark}

\begin{remark}[Conjecture~1.2 and exotic smooth structures]
Axiom~\ref{ax:selfconsistency} includes the conjecture that multiple
fine-grained configurations are compatible with the same coarse-grained
solution. Theorem~\ref{thm:exotic} is the mathematical realization of this
conjecture at Level~1: a single PL 4-manifold (the coarse-grained description)
admits uncountably many smooth structures (the fine-grained configurations).
This is not an analogy. It is, at the smooth-structure level, the precise
content of the conjecture.
\end{remark}

%----------------------------------------------------------------------
\subsection{Level 2: The metric fixed point}
\label{sec:level2}
%----------------------------------------------------------------------

Given a smooth structure $\sigma \in \mathrm{Sm}(\mathcal{M})$, the
self-consistency constraint at the metric level is:

\begin{equation}
G_{\mu\nu}[g] = 8\pi G\,\langle\hat{T}_{\mu\nu}\rangle_{g,\mathrm{ren}}
\label{eq:sce}
\end{equation}

This is the semiclassical Einstein equation, which is a fixed-point condition
on the metric $g$: the geometry that hosts the quantum state must be sourced by
that state. The companion paper~\cite{fixed_point} establishes existence of
solutions at three levels of rigor:

\begin{enumerate}

\item \emph{Exactly}, in the cosmological case via the Starobinsky
(1980) trace anomaly solution, with $H_0^2 = 180\pi/G|a_2|$ (Rigorous).

\item \emph{Perturbatively}, near any classical solution, via a Banach
contraction with constant $\kappa \sim (m/M_P)^2 \ll 1$ for massive fields
(Sketch; the $H^s$ operator-norm bound is open). The extension to massless
fields (photon, gluon) remains open: the spectral gap that controls the
infrared behavior of the response kernel vanishes at $m = 0$.

\item \emph{Conditionally non-perturbatively}, via the Schauder fixed-point
theorem on globally hyperbolic spacetimes with compact Cauchy surfaces, under
six stated assumptions of which four are established and two remain open
(Sketch).

\end{enumerate}

The Planck scale emerges as the natural validity boundary: when $R \sim
\ell_P^{-2}$, the contraction constant $\kappa \to 1$, the perturbative
expansion breaks down, and the full smooth-structure measure $\mu^*$ becomes
relevant.

\begin{remark}[Level~2 is a coarse-graining of Level~1]
The companion paper's self-consistency map $\mathcal{F}: g \mapsto g'$ operates
on single metrics. This is the special case of the smooth-structure map
$\mathcal{F}$ of Definition~\ref{def:mu} in which $\mu$ is a delta measure
concentrated on a single smooth structure, and the smooth structure, together
with the Level~2 fixed-point condition, is conjectured to determine a unique
metric. The companion paper establishes that, in this limiting case,
the fixed point exists. The full conjecture is that the fixed point of $\mathcal{F}$
on the space of measures over smooth structures exists and is unique.
\end{remark}

%----------------------------------------------------------------------
\subsection{Level 3: Effective quantum mechanics}
\label{sec:level3}
%----------------------------------------------------------------------

\begin{conjecture}[Emergence of effective quantum mechanics]
\emph{(Conjecture.)}
\label{conj:qm}
Let $\mathcal{R} \subset \mathcal{M}$ be a subregion (a ``laboratory'' in the
block). Let $\mu^*$ be the self-consistency measure on
$\mathrm{Sm}(\mathcal{M})$. The marginal of $\mu^*$ over the smooth structures
of the complement $\mathcal{M} \setminus \mathcal{R}$, restricted to
$\mathcal{R}$, defines an effective measure $\mu^*_\mathcal{R}$ on the local
configuration space of $\mathcal{R}$.

When $\mathcal{R}$ is much larger than the Planck scale but much smaller than
the cosmological scale, and gravity within $\mathcal{R}$ is weak, $\mu^*_\mathcal{R}$
is approximated by a density matrix $\rho_\mathcal{R}$ on the Hilbert space
$\mathcal{H}_\mathcal{R} = L^2(\mathrm{Config}(\mathcal{R}), \mu^*_\mathcal{R})$.
The quantum mechanical statistics seen by subsystems within $\mathcal{R}$ are
a consequence of this marginalization.
\end{conjecture}

\begin{remark}[Why marginalization produces quantum statistics]
The argument for Conjecture~\ref{conj:qm} parallels the derivation of the
canonical ensemble from microcanonical statistics: integrating out the
environment imposes linearity on the subsystem. In the present context, the
consistency constraints on the full block, restricted to $\mathcal{R}$, become
approximately linear in the degrees of freedom of $\mathcal{R}$ when the
complement is integrated over, because the dominant contribution from the
complement is the mean-field background, with fluctuations suppressed by
$V_\mathcal{R}/V_{\mathrm{total}}$. The effective measure on $\mathcal{R}$ is
thus a density matrix. The analogy to statistical mechanics is structural---in
both cases, integrating out a large environment produces an approximately linear
effective theory for the subsystem---but the mechanism differs. In statistical
mechanics, the linearity of the canonical ensemble follows from the specific
exponential (Boltzmann) form of the microcanonical distribution. The
self-consistency measure $\mu^*$ has no \textit{a priori} reason to take
exponential form; whether the marginalization of $\mu^*$ produces the same
linearization is the content of Conjecture~\ref{conj:qm}, not a consequence of
the analogy. This argument is a conjecture; it requires a proof that
the self-consistency constraints, when restricted and marginalized, linearize
in the subsystem's degrees of freedom.
\end{remark}

\begin{remark}[The Hilbert space is not fundamental]
Conjecture~\ref{conj:qm} implies that the Hilbert space $\mathcal{H}_\mathcal{R}$
is:
\begin{enumerate}
\item \emph{Derived}, not postulated: it is the $L^2$ space of the marginal measure.
\item \emph{Subsystem-dependent}: different choices of $\mathcal{R}$ give different
Hilbert spaces. There is no single Hilbert space of the universe.
\item \emph{Approximate}: the density matrix approximation is valid only in the
sub-Planckian, weak-gravity regime.
\end{enumerate}
This is consistent with the known difficulties of defining a universal Hilbert
space in quantum gravity, and with the Page--Wootters result~\cite{page_wootters}
that relational time emerges from a globally static state via conditioning on a
subsystem.
\end{remark}

%----------------------------------------------------------------------
\subsection{The interlocking constraint}
\label{sec:interlock}
%----------------------------------------------------------------------

The four levels are not independent. The constraint is that the self-consistency
chain must close:

\begin{equation}
\underbrace{Q^*}_{\text{Level 0}}
\xrightarrow{\text{Freedman/Donaldson}}
\underbrace{\mathrm{Sm}(\mathcal{M})}_{\text{Level 1}}
\xrightarrow{\mu^* = \mathcal{F}(\mu^*)}
\underbrace{g^*}_{\text{Level 2}}
\xrightarrow{\text{Einstein}}
\underbrace{\langle\hat{T}_{\mu\nu}\rangle}_{\text{Level 2}}
\xrightarrow{\text{marginalize}}
\underbrace{\rho_\mathcal{R}}_{\text{Level 3}}
\xrightarrow{\text{trace back}}
Q^*
\label{eq:chain}
\end{equation}

The arrows represent logical constraints, not temporal sequences. A block
satisfying the hierarchy satisfies all constraints simultaneously. The chain
\eqref{eq:chain} is closed: the matter content derived from the effective
quantum state $\rho_\mathcal{R}$ at Level~3 must be consistent with the
intersection form $Q^*$ at Level~0 that initiated the chain.

\begin{remark}[The chain is an epistemic decomposition, not ontological stages]
The levels in \eqref{eq:chain} are introduced for human comprehensibility.
They correspond to the precision at which we can currently do mathematics. The
block does not ``proceed'' through these levels. It satisfies the constraint.
The iteration suggested by Conjecture~1.1 of Axiom~\ref{ax:selfconsistency}
is our procedure for computing the fixed point, not a description of how the
universe came to be.
\end{remark}

\begin{remark}[The ``trace back'' arrow]
\label{rem:traceback}
The final arrow in~\eqref{eq:chain}, $\rho_\mathcal{R} \to Q^*$, is the least
developed constraint in the chain. The claim is that the effective quantum state
of subsystems, aggregated over all subregions, must be consistent with the
global intersection form that initiated the chain. Formally, this requires:
\begin{enumerate}
\item[(i)] A reconstruction theorem showing that local density matrices
determine a global metric (an inverse problem).
\item[(ii)] A topological rigidity result showing that the global metric
determines the intersection form (known for some classes of 4-manifolds).
\item[(iii)] A self-consistency proof that the resulting $Q^*$ is the same one
that initiated the chain.
\end{enumerate}
None of (i)--(iii) is established. The ``trace back'' arrow is a conjecture
about a conjecture, and making it precise is prerequisite to any claim that the
chain~\eqref{eq:chain} actually closes.
\end{remark}

%======================================================================
\section{The dimensionality argument}
\label{sec:dimensionality}
%======================================================================

The hierarchy of Section~\ref{sec:hierarchy} presupposes a 4-manifold. Here
we argue that the four-dimensionality is not an input but a consequence: four
is the unique dimension where the constraint structure of
Definition~\ref{def:hierarchy} is nontrivial.

\begin{table}[h]
\centering
\begin{tabular}{lccc}
\toprule
Dimension & Smooth structures on $\mathbb{R}^n$ & $h$-cobordism theorem & Level~0 richness \\
\midrule
$n < 4$ & Unique & Holds (Smale, $n \geq 6$; special cases $n = 5$) & Insufficient \\
$n = 4$ & Uncountably many & \textbf{Fails} (Donaldson 1987) & Rich \\
$n > 4$ & Unique ($n \geq 5$) & Holds & Insufficient \\
\bottomrule
\end{tabular}
\caption{Topology by dimension. The failure of the $h$-cobordism theorem in
dimension four is the root of the exotic structure richness.}
\label{tab:dimensions}
\end{table}

\begin{proposition}[Four-dimensional uniqueness]
\emph{(Sketch.)}
\label{prop:4d}
The self-consistency hierarchy of Definition~\ref{def:hierarchy} is trivial in
every dimension other than four:
\begin{enumerate}
\item For $n \neq 4$: $\mathrm{Sm}(\mathcal{M})$ is a finite set (generically
a single point for $\mathbb{R}^n$). The Level~1 measure $\mu^*$ is a finite
sum, not an integral over a continuous moduli space. The effective quantum
mechanics of Conjecture~\ref{conj:qm} collapses to classical statistics.

\item For $n = 4$: $\mathrm{Sm}(\mathcal{M})$ is uncountably infinite.
The Level~1 measure is a genuine integral over a continuous moduli space.
The Level~0 constraint is nontrivial because the Freedman--Donaldson gap is
maximal: many topological manifolds exist that admit no smooth structure at all.
\end{enumerate}
\end{proposition}

\noindent\emph{Sketch.} Point (1) uses Theorem~\ref{thm:exotic}: smooth
structures on $\mathbb{R}^n$ are unique for $n \neq 4$. The same conclusion
holds for compact manifolds in dimensions $n \geq 5$ (Smale's $h$-cobordism
theorem; see~\cite{smale}), and for $n \leq 3$ by classical surface and
3-manifold theory. For $n = 4$, Theorems~\ref{thm:freedman}
and~\ref{thm:donaldson} establish the Freedman--Donaldson gap, and
Theorem~\ref{thm:exotic} gives the uncountable smooth structures.
Point (2) is the content of these theorems. The identification of the
uncountable smooth structure moduli space as the source of quantum statistics
(Conjecture~\ref{conj:qm}) is a conjecture, not a theorem. \hfill$\square$

\begin{remark}
Proposition~\ref{prop:4d} does not \emph{prove} that the universe is four-dimensional.
It establishes that the framework proposed here is nontrivial only in dimension
four. If the framework is correct, four-dimensionality follows. Whether the
framework is correct is an open question.
\end{remark}

%======================================================================
\section{Connection to existing formulations}
\label{sec:connection}
%======================================================================

\subsection{The Wheeler--DeWitt equation as a Level-2 projection}

The Wheeler--DeWitt equation $\hat{H}\Psi[\gamma_{ij}] = 0$~\cite{dewitt} is
the canonical quantization of the metric-level constraint. In the hierarchy of
Section~\ref{sec:hierarchy}, it is an approximation to the Level~2 fixed-point
condition, obtained by:

\begin{enumerate}
\item Concentrating the smooth-structure measure $\mu^*$ on a single smooth
structure (ignoring the Level~1 sum over exotic structures).
\item Performing a $3+1$ decomposition of the metric (introducing a preferred
foliation).
\item Replacing the Level~2 classical fixed-point condition with its quantum
analogue via Dirac quantization.
\end{enumerate}

Each step introduces additional structure not present in the hierarchy. Step (2)
is the primary source of Axiom~\ref{ax:timeless} violation~\cite{kuchar}.
Steps (1) and (3) are simplifications
that are valid in the sub-Planckian, single-smooth-structure regime.

The Wheeler--DeWitt equation is thus the projection of the full hierarchy onto
a single smooth structure and a single spatial foliation. It contains the
correct physics at leading order in $\kappa \sim (m/M_P)^2$, but not the
global, foliation-independent constraint structure of the full hierarchy.

\subsection{The semiclassical Einstein equation as Level-2}

The companion paper~\cite{fixed_point} studies eq.~\eqref{eq:sce} as a
fixed-point condition in its own right. Within the present hierarchy, this is
the Level~2 constraint evaluated on a single smooth structure. The Banach
contraction result ($\kappa \ll 1$ for sub-Planckian curvatures) is the
perturbative stability of the fixed point at this level. The Schauder existence
result (conditional on two open conjectures) is non-perturbative existence.

The hierarchy extends this program: Level~2 existence (established with caveats
in the companion paper) is a necessary condition for the full hierarchy to have
a fixed point, but not sufficient. Level~1 requires additionally that the
measure over smooth structures has a fixed point. Level~0 requires that the
intersection form is self-consistent with the matter content.

\subsection{The Page--Wootters mechanism at Level~3}

The Page--Wootters mechanism~\cite{page_wootters}---apparent time evolution
of a subsystem recovered from a globally static state via conditioning on a
``clock'' subsystem---is a Level~3 phenomenon in the present hierarchy. The
globally static state is the smooth-structure measure $\mu^*$; the
conditioning on a subsystem is the marginalization of Conjecture~\ref{conj:qm}.
The apparent time evolution is a pattern in the correlations between
$\mathcal{R}$ and its complement encoded in $\mu^*$.

The observer-dependence identified by Kucha\v{r}~\cite{kuchar} and confirmed
by Marletto and Vedral~\cite{marletto_vedral}---different choices of ``clock''
subsystem give different dynamics---is, in the present hierarchy, a manifestation
of the fact that different choices of $\mathcal{R}$ give different marginal
measures $\mu^*_\mathcal{R}$ and hence different effective Hilbert spaces
$\mathcal{H}_\mathcal{R}$. The observer-dependence is a feature, not a bug:
it reflects the fact that the Hilbert space is emergent and subsystem-relative,
not fundamental.

%======================================================================
\section{The metric-to-signature problem}
\label{sec:signature}
%======================================================================

The most critical open problem in the hierarchy concerns the emergence of
Lorentzian metric signature. We identify it precisely here because it is
potentially fatal to the topological approach.

The intersection form $Q_\mathcal{M}$ has \emph{type} $(b^+, b^-)$, where
$b^\pm$ are the dimensions of the positive- and negative-definite subspaces,
and \emph{signature} $\sigma(\mathcal{M}) = b^+ - b^-$. The Hirzebruch
signature theorem~\cite{hirzebruch} gives $\sigma(\mathcal{M}) =
\frac{1}{3}p_1(\mathcal{M})$, connecting the intersection form to the first
Pontryagin class. The intersection form type $(b^+, b^-)$ is a global invariant
of $H_2(\mathcal{M};\mathbb{Z})$. It must be carefully distinguished from the
\emph{metric signature}---the sign pattern $(-,+,+,+)$ of the tangent-space
metric---which is a pointwise property of a different mathematical object.

The conjectured mechanism for Lorentzian signature emergence is:

\begin{conjecture}[Signature emergence]
\emph{(Conjecture.)}
\label{conj:signature}
An indefinite intersection form $Q_\mathcal{M}$ of type $(b^+, b^-)$ with
$b^+ = 1$---one positive eigenvalue and $b^-$ negative eigenvalues---requires
that any smooth metric compatible with the self-consistency constraint of
Level~2 has Lorentzian metric signature $(-,+,+,+)$. No known theorem connects
the intersection form type to the metric signature; bridging this gap between
$H_2$ and the tangent bundle is the central mathematical obstacle.
\end{conjecture}

We emphasize that Conjecture~\ref{conj:signature} is the weakest link in the
hierarchy. The intersection form is a global invariant on $H_2$; the metric
signature is a local property of the tangent bundle. Connecting global topology
to local geometry requires a theorem that, to our knowledge, does not currently
exist in the mathematical literature.

The direction in which such a theorem might go: an indefinite form of type
$(1, b^-)$ has $b^+ = 1$, which is \emph{suggestive} of a single timelike
direction, but this is an analogy between second homology and the tangent
bundle that currently has no mathematical justification. Definite forms
correspond to Riemannian (positive-definite) metrics; indefinite forms with
different types correspond to other pseudo-Riemannian structures. The specific
$b^+ = 1$ case is singled out by requiring:
\begin{enumerate}
\item The self-consistent manifold must admit a global timelike direction (for
the notion of ``subsystem'' in Level~3 to be well-defined).
\item The signature must be consistent with the causal structure required by the
Level~2 Einstein equation.
\end{enumerate}
These requirements, formalized, might yield Conjecture~\ref{conj:signature}.
The proof remains open.

%======================================================================
\section{Open problems}
\label{sec:open}
%======================================================================

We list four priority open problems, in order of importance:

\begin{enumerate}

\item \textbf{The signature bridge (Section~\ref{sec:signature}).}
Prove or disprove Conjecture~\ref{conj:signature}: that an indefinite
intersection form of type $(1,b^-)$ with $b^+ = 1$ implies Lorentzian metric
signature $(-,+,+,+)$ for the self-consistent metric. This requires bridging
the gap between a global invariant on $H_2$ and a local property of the tangent
bundle---a specific mathematical question in four-dimensional topology and
Lorentzian geometry. A proof would
validate the topological approach to signature emergence. A counterexample would
require a different mechanism for signature emergence or a modification of the
Level~0 constraint.

\item \textbf{A natural measure on smooth structures.}
Define a measure $\mu$ on $\mathrm{Sm}(\mathcal{M})$ for a compact
4-manifold $\mathcal{M}$ that is:
\begin{enumerate}
\item Defined without reference to a metric (to avoid the circularity of the
standard gravitational measure problem).
\item Compatible with the self-consistency condition $\mathcal{F}(\mu) = \mu$.
\item Invariant under the orientation-preserving diffeomorphism group.
\end{enumerate}
Donaldson--Witten theory and Seiberg--Witten invariants provide natural weights
on smooth structures; assembling these into a well-defined measure on the full
moduli space $\mathrm{Sm}(\mathcal{M})$ is the open problem.

\item \textbf{The marginalization conjecture (Conjecture~\ref{conj:qm}).}
Prove that the marginal of a self-consistent measure $\mu^*$ on
$\mathrm{Sm}(\mathcal{M})$, restricted to a subregion $\mathcal{R}$ in
the sub-Planckian, weak-gravity regime, is approximated by a density matrix
on a Hilbert space. Even a toy model---a finite-dimensional analogue of the
full measure theory---demonstrating that marginalization of a self-consistent
measure produces quantum statistics would be significant.

\item \textbf{Connecting the topological constraint to the Banach contraction.}
The companion paper~\cite{fixed_point} establishes a Banach contraction at
Level~2 with constant $\kappa \sim (m/M_P)^2$. Is there an analogous
contraction estimate at Level~1 (on the space of smooth-structure measures)?
Such an estimate would extend the companion paper's uniqueness and convergence
results from single metrics to the full measure on smooth structures.

\end{enumerate}

%======================================================================
\begin{thebibliography}{99}

\bibitem{dewitt} B.~S.~DeWitt, Quantum theory of gravity.~I.~The canonical
theory, \textit{Phys.\ Rev.}\ \textbf{160}, 1113 (1967).

\bibitem{wheeler} J.~A.~Wheeler, Superspace and the nature of quantum
geometrodynamics, in \textit{Battelle Rencontres: 1967 Lectures in
Mathematics and Physics}, eds.\ C.~DeWitt and J.~A.~Wheeler (Benjamin,
New York, 1968).

\bibitem{freedman} M.~H.~Freedman, The topology of four-dimensional manifolds,
\textit{J.\ Differential Geometry}\ \textbf{17}, 357--453 (1982).

\bibitem{donaldson} S.~K.~Donaldson, An application of gauge theory to
four-dimensional topology, \textit{J.\ Differential Geometry}\ \textbf{18},
279--315 (1983).

\bibitem{taubes} C.~H.~Taubes, Gauge theory on asymptotically periodic
$4$-manifolds, \textit{J.\ Differential Geometry}\ \textbf{25}, 363--430 (1987).

\bibitem{smale} S.~Smale, Generalized Poincar\'e's conjecture in dimensions
greater than four, \textit{Ann.\ Math.}\ \textbf{74}, 391--406 (1961).

\bibitem{hirzebruch} F.~Hirzebruch, \textit{Topological Methods in Algebraic
Geometry}, 3rd ed., Springer, Berlin (1966).

\bibitem{page_wootters} D.~N.~Page and W.~K.~Wootters, Evolution without
evolution: Dynamics described by stationary observables,
\textit{Phys.\ Rev.\ D}\ \textbf{27}, 2885 (1983).

\bibitem{kuchar} K.~Kucha\v{r}, Time and interpretations of quantum
gravity, in \textit{Proceedings of the 4th Canadian Conference on General
Relativity and Relativistic Astrophysics}, eds.\ G.~Kunstatter,
D.~Vincent, and J.~Williams (World Scientific, Singapore, 1992).

\bibitem{marletto_vedral} C.~Marletto and V.~Vedral, Evolution without
evolution and without ambiguities, \textit{Phys.\ Rev.\ D}\ \textbf{95},
043510 (2017).

\bibitem{fixed_point} W.~Regelmann and Claude (Anthropic), Fixed-point
existence of self-consistent semiclassical gravity (2026). Companion paper.

\bibitem{halliwell} J.~J.~Halliwell, Derivation of the Wheeler--DeWitt
equation from a path integral for minisuperspace models,
\textit{Phys.\ Rev.\ D}\ \textbf{38}, 2468 (1988).

\bibitem{banks} T.~Banks, TCP, quantum gravity, the cosmological constant
and all that\ldots, \textit{Nucl.\ Phys.\ B}\ \textbf{249}, 332 (1985).

\end{thebibliography}

\end{document}
